\chapter{System Architecture and UML Modeling}

\section{Architectural Vision}
The SmartHome platform follows a layered event-driven architecture combining:
\begin{itemize}[leftmargin=1.2cm]
  \item API-oriented business services (REST).
  \item Event-oriented machine communication (MQTT).
  \item Real-time user synchronization (Socket.IO).
  \item Persistent state and history (MongoDB).
\end{itemize}

This architecture separates concerns while allowing fast propagation of state changes across subsystems.

\section{High-Level Context Diagram}
\begin{figure}[H]
\centering
\resizebox{\textwidth}{!}{%
\begin{tikzpicture}[
  node distance=1.2cm and 1.6cm,
  actor/.style={draw, rounded corners, fill=blue!10, minimum width=3.2cm, minimum height=0.9cm, align=center},
  system/.style={draw, rounded corners, fill=green!12, minimum width=5.0cm, minimum height=1.2cm, align=center},
  datastore/.style={draw, cylinder, shape border rotate=90, aspect=0.2, minimum height=1.2cm, minimum width=2.0cm, fill=orange!15, align=center}
]
\node[actor] (agency) {Agency User};
\node[actor, below=of agency] (admin) {House Admin};
\node[actor, below=of admin] (resident) {Resident};

\node[system, right=3.2cm of admin] (platform) {SmartHome Platform\\(API + Realtime + Broker)};
\node[datastore, right=3.6cm of platform] (db) {MongoDB};
\node[actor, below=1.6cm of platform] (devices) {IoT Devices};

\draw[-{Latex[length=3mm]}] (agency) -- node[above,sloped]{Manage units, requests} (platform);
\draw[-{Latex[length=3mm]}] (admin) -- node[above,sloped]{Control home, review permissions} (platform);
\draw[-{Latex[length=3mm]}] (resident) -- node[below,sloped]{Use allowed features} (platform);
\draw[-{Latex[length=3mm]}] (platform) -- node[above]{Persist/read} (db);
\draw[-{Latex[length=3mm]}] (devices) -- node[right]{MQTT telemetry / commands} (platform);
\end{tikzpicture}
}
\caption{System context diagram}
\end{figure}

\section{Component Diagram}
\begin{figure}[H]
\centering
\resizebox{\textwidth}{!}{%
\begin{tikzpicture}[
  comp/.style={draw, rounded corners, fill=cyan!10, minimum width=3.4cm, minimum height=1.0cm, align=center},
  arr/.style={-{Latex[length=3mm]}}
]
\node[comp] (frontend) at (0,0) {React Frontend};
\node[comp] (api) at (5,0) {Express API Layer};
\node[comp] (socket) at (10,0) {Socket.IO Service};

\node[comp] (auth) at (5,-2) {Auth and RBAC Middleware};
\node[comp] (domain) at (5,-4) {Domain Routes and Services};
\node[comp] (broker) at (10,-2) {Aedes MQTT Broker};
\node[comp] (gas) at (10,-4) {Gas Monitor Service};
\node[comp] (mongo) at (5,-6) {MongoDB and Models};

\draw[arr] (frontend) -- node[above]{REST} (api);
\draw[arr] (frontend) -- node[above]{Realtime events} (socket);
\draw[arr] (api) -- (auth);
\draw[arr] (auth) -- (domain);
\draw[arr] (domain) -- (mongo);
\draw[arr] (broker) -- node[right]{publish/subscribe} (gas);
\draw[arr] (gas) -- (mongo);
\draw[arr] (gas) -- node[left]{emit alerts} (socket);
\draw[arr] (domain) -- (broker);
\draw[arr] (api) -- (socket);
\end{tikzpicture}
}
\caption{Logical component architecture}
\end{figure}

\section{Deployment Diagram}
\begin{figure}[H]
\centering
\resizebox{\textwidth}{!}{%
\begin{tikzpicture}[
  box/.style={draw, rounded corners, fill=gray!10, minimum width=4.4cm, minimum height=1.2cm, align=center},
  arr/.style={-{Latex[length=3mm]}}
]
\node[box] (browser) at (0,0) {Client Browser\\(React SPA)};
\node[box] (node) at (6,0) {Node.js Runtime\\Express + Socket.IO + Aedes};
\node[box] (mongo) at (12,0) {MongoDB Server};
\node[box] (iot) at (6,-2.2) {IoT Devices / Simulators\\MQTT TCP / MQTT over WS};
\node[box] (mail) at (12,-2.2) {Email Service\\(verification/reset)};

\draw[arr] (browser) -- node[above]{HTTPS/HTTP + WS} (node);
\draw[arr] (node) -- node[above]{Mongoose} (mongo);
\draw[arr] (iot) -- node[left]{MQTT topics} (node);
\draw[arr] (node) -- node[right]{SMTP/API} (mail);
\end{tikzpicture}
}
\caption{Deployment-level topology}
\end{figure}

\section{Use Case Diagram Narrative}
The main use-case clusters are:
\begin{itemize}[leftmargin=1.2cm]
  \item Agency governance: manage units, review owner requests, monitor analytics.
  \item Admin operations: manage users and devices, approve permissions, activate scenes.
  \item Resident interactions: controlled dashboard usage, permission requests, profile actions.
  \item Automated safety loop: ingest sensor values, compute emergency state, propagate alerts.
\end{itemize}

\section{Sequence Diagram: Resident Permission Request}
\begin{figure}[H]
\centering
\resizebox{\textwidth}{!}{%
\begin{tikzpicture}[
  lifeline/.style={draw, minimum width=2.6cm, minimum height=0.8cm, fill=blue!8, align=center},
  arr/.style={-{Latex[length=2.5mm]}}
]
\node[lifeline] (res) at (0,0) {Resident UI};
\node[lifeline] (api) at (4,0) {Permissions API};
\node[lifeline] (db) at (8,0) {MongoDB};
\node[lifeline] (sock) at (12,0) {Socket.IO};
\node[lifeline] (adm) at (16,0) {Admin UI};

\draw[dashed] (res) -- (0,-5.2);
\draw[dashed] (api) -- (4,-5.2);
\draw[dashed] (db) -- (8,-5.2);
\draw[dashed] (sock) -- (12,-5.2);
\draw[dashed] (adm) -- (16,-5.2);

\draw[arr] (0,-1) -- node[above]{POST request} (4,-1);
\draw[arr] (4,-2) -- node[above]{create doc} (8,-2);
\draw[arr] (8,-2.6) -- node[below]{created} (4,-2.6);
\draw[arr] (4,-3.4) -- node[above]{emit event} (12,-3.4);
\draw[arr] (12,-4.2) -- node[above]{requestCreated} (16,-4.2);
\end{tikzpicture}
}
\caption{Permission request event propagation}
\end{figure}

\section{Sequence Diagram: Gas Emergency}
\begin{figure}[H]
\centering
\resizebox{\textwidth}{!}{%
\begin{tikzpicture}[
  lifeline/.style={draw, minimum width=2.4cm, minimum height=0.8cm, fill=red!8, align=center},
  arr/.style={-{Latex[length=2.5mm]}}
]
\node[lifeline] (sensor) at (0,0) {Gas Sensor};
\node[lifeline] (broker) at (3.8,0) {MQTT Broker};
\node[lifeline] (monitor) at (7.6,0) {Gas Monitor};
\node[lifeline] (db) at (11.4,0) {HouseState DB};
\node[lifeline] (sock) at (15.2,0) {Socket.IO};
\node[lifeline] (ui) at (19,0) {Dashboards};

\foreach \x in {0,3.8,7.6,11.4,15.2,19} {\draw[dashed] (\x,-0.4) -- (\x,-6.2);} 

\draw[arr] (0,-1.0) -- node[above]{publish sensors/H1234/gas} (3.8,-1.0);
\draw[arr] (3.8,-1.8) -- node[above]{event packet} (7.6,-1.8);
\draw[arr] (7.6,-2.6) -- node[above]{save level + emergencyMode} (11.4,-2.6);
\draw[arr] (7.6,-3.4) -- node[above]{publish alarm + valve command} (3.8,-3.4);
\draw[arr] (7.6,-4.2) -- node[above]{emit house:gas-emergency} (15.2,-4.2);
\draw[arr] (15.2,-5.0) -- node[above]{real-time alert} (19,-5.0);
\end{tikzpicture}
}
\caption{Gas emergency sequence}
\end{figure}

\section{Activity Diagram Narrative: User Onboarding}
User onboarding branch logic:
\begin{enumerate}[leftmargin=1.2cm]
  \item User authenticates.
  \item If house code already exists in profile, redirect to dashboard.
  \item Else choose between join existing unit and create solo unit.
  \item Persist onboarding result and assign role capabilities.
  \item Navigate to role-appropriate dashboard.
\end{enumerate}

\section{State Diagram Narrative: House Safety State}
House safety states can be conceptualized as:
\begin{itemize}[leftmargin=1.2cm]
  \item Normal: gas level below threshold, emergency disabled.
  \item Alert Triggered: gas level exceeds threshold and transition detected.
  \item Emergency Active: emergency mode true, valve closed, alarm command retained.
  \item Recovery: value returns below threshold or admin clear action executed.
\end{itemize}

\section{Architectural Quality Attributes}
The architecture optimizes for:
\begin{itemize}[leftmargin=1.2cm]
  \item Modularity: independent route and service modules.
  \item Responsiveness: push events to clients instead of polling.
  \item Evolvability: model and route extensions with localized impact.
  \item Operational clarity: deterministic handling for safety events.
\end{itemize}

\section{Design Constraints and Evolution Strategy}
The architecture currently co-hosts API, sockets, and broker in one process. This is acceptable for project and early deployment stages. A future production strategy can separate broker, API, and realtime gateway into independent services while preserving contracts defined in this chapter.
